%% ma: L10-B13-0050
%% lop: 10
%% chuong: 5
%% bai: 13
%% dang: 10.13.6
%% loai: DS
%% muc: [NB, NB, NB, TH]
%% dap_an: "ĐSSĐ"
%% nguon: "chuyen-de-cac-so-dac-trung-cua-mau-so-lieu-khong-ghep-nhom-toan-10-chuong-trinh-moi.pdf (GV. Phan Nhật Linh) (P02), Bài 02 Dạng 1 Phần II Câu 5"
%% kien_thuc: "Lớp 10 Bài 13 (số trung bình)"
%% trang_thai: chinh-thuc
%% ghi_chu: "Tự duyệt theo yêu cầu giáo viên 22/10/2026."
\begin{ex}
Tình hình thu hoạch lúa vụ mùa năm 2022 của ba hợp tác xã ở một địa phương như sau:
\begin{center}
\begin{tabular}{|c|c|c|}\hline
Hợp tác xã & Năng suất lúa (tạ/ha) & Diện tích trồng lúa (ha)\\\hline
A & $40$ & $150$\\\hline
B & $38$ & $130$\\\hline
C & $36$ & $120$\\\hline
\end{tabular}
\end{center}
Xét tính đúng sai của các khẳng định sau:
\choiceTF{\True Sản lượng lúa của hợp tác xã A là $6000$ tạ}{Sản lượng lúa của hợp tác xã B là $4950$ tạ}{Sản lượng lúa của hợp tác xã C là $4120$ tạ}{\True Năng suất lúa trung bình của toàn bộ ba hợp tác xã là $38{,}15$ tạ/ha}
\loigiai{a) Đúng: $40\cdot150=6000$.

b) Sai: $38\cdot130=4940$.

c) Sai: $36\cdot120=4320$.

d) Đúng: $\dfrac{6000+4940+4320}{150+130+120}=\dfrac{15\,260}{400}=38{,}15$ (tạ/ha).}
\end{ex}
